\documentclass[11pt]{article}

\usepackage[T1]{fontenc}
\usepackage[utf8]{inputenc}
\usepackage{amsmath,amssymb,amsthm,mathtools}
\usepackage{newtxtext,newtxmath}
\usepackage[margin=1.08in]{geometry}
\usepackage{microtype}
\usepackage{enumitem}
\usepackage{xcolor}
\usepackage{hyperref}
\usepackage[nameinlink,noabbrev]{cleveref}
\usepackage{xurl}

\definecolor{linkblue}{HTML}{174A7E}
\definecolor{muted}{HTML}{5D6670}
\hypersetup{
  colorlinks=true,
  linkcolor=linkblue,
  citecolor=linkblue,
  urlcolor=linkblue,
  pdftitle={David Roe's Candidate Presentation for the Absolute Galois Group of Q2},
  pdfsubject={A verification by the finite-quotient method},
  pdfkeywords={absolute Galois group, Q2, profinite presentation, Demushkin group}
}

\setlist[itemize]{leftmargin=1.6em,itemsep=0.25em,topsep=0.4em}
\setlist[enumerate]{leftmargin=1.8em,itemsep=0.25em,topsep=0.4em}
\allowdisplaybreaks
\numberwithin{equation}{section}

\newtheorem{theorem}{Theorem}[section]
\newtheorem{proposition}[theorem]{Proposition}
\newtheorem{lemma}[theorem]{Lemma}
\newtheorem{corollary}[theorem]{Corollary}
\theoremstyle{definition}
\newtheorem{definition}[theorem]{Definition}
\theoremstyle{remark}
\newtheorem{remark}[theorem]{Remark}

\newcommand{\Q}{\mathbf Q}
\newcommand{\Z}{\mathbf Z}
\newcommand{\F}{\mathbf F}
\newcommand{\GQ}{G_{\Q_2}}
\newcommand{\GQpro}{G_{\Q_2}(2)}
\newcommand{\Gal}{\operatorname{Gal}}
\newcommand{\Sur}{\operatorname{Sur}}
\newcommand{\Aut}{\operatorname{Aut}}
\newcommand{\im}{\operatorname{im}}
\newcommand{\tr}{\operatorname{tr}}
\newcommand{\ab}{\mathrm{ab}}
\newcommand{\ur}{\mathrm{ur}}
\newcommand{\rt}{r_{\mathrm t}}
\newcommand{\rR}{r_R}
\newcommand{\GammaR}{\Gamma_R}
\newcommand{\TR}{\mathcal T}
\newcommand{\DR}{D_R}
\newcommand{\Dzero}{D_0}
\newcommand{\nuR}{\nu_R}
\newcommand{\chiR}{\chi_R}
\newcommand{\chiZero}{\chi_0}
\newcommand{\omegaTwo}{\omega_2}
\newcommand{\Zassenhaus}[1]{\mathcal D_{#1}}
\newcommand{\angles}[1]{\left\langle #1\right\rangle}
\newcommand{\normalclosure}[1]{\overline{\left\langle\!\left\langle #1
  \right\rangle\!\right\rangle}}

\title{\vspace{-2.5em}
  \textbf{David Roe's Candidate Presentation for the\\
  Absolute Galois Group of \(\Q_2\)}\\[0.45em]
  \large A verification by the finite-quotient method}
\author{}
\date{July 24, 2026}

\begin{document}
\maketitle
\vspace{-2.4em}

\begin{abstract}
We verify the candidate presentation for \(G_{\Q_2}\) listed as David Roe's
candidate on the presentation website \cite{Presentations}.  The verification
is relative to the general local-cohomological and finite-target results of
Roe--Turturean \cite{RT}: we prove that replacing their candidate word by
Roe's word preserves every source-side input to their boundary-framed
finite-quotient induction.  The main new point is the integral marking of the
maximal pro-\(2\) quotient.  Its canonical Demushkin orientation is computed
directly from Labute's crossed-derivation characterization.  The evaluated
Fox matrix is then the matrix of \cite[\S5]{RT} with the two wild columns
interchanged, and the Fox--Heisenberg pairing and determinant quadratic form
are unchanged after the same interchange.  The induction of
\cite[\S\S7--10]{RT} consequently gives equality of epimorphism counts to
every finite group and hence the asserted profinite isomorphism.
\end{abstract}

\begin{center}
\small\color{muted}
\textbf{Logical status.}
All calculations depending on Roe's relator are proved below.
The local Tate-duality, Gauss-sign, Fourier-inversion, target-induction, and
profinite-reconstruction theorems are invoked from \cite{RT}; they are not
reproved here.
\end{center}

\tableofcontents

\section{The candidate and the replacement theorem}

We use the conventions
\[
  g^h=h^{-1}gh,\qquad [g,h]=g^{-1}h^{-1}gh.
\]
Let
\[
  \omegaTwo=(1\text{ in }\Z_2,\ 0\text{ in }\Z_\ell
  \text{ for every odd prime }\ell)\in\widehat{\Z}.
\]
Thus \(\omegaTwo^2=\omegaTwo\), and exponentiation by \(\omegaTwo\) extracts
the \(2\)-primary component of the procyclic group generated by an element.
Put
\begin{equation}\label{eq:defwords}
  \sigma_2=\sigma^{\omegaTwo},\qquad
  a=(x_0^{-3}\tau)^{\omegaTwo},\qquad
  y_1=x_1^{\sigma_2},\qquad
  c=[x_1,y_1].
\end{equation}
The two relators are
\begin{equation}\label{eq:relators}
  \rt=(\tau^\sigma)^{-1}\tau^2,\qquad
  \rR=(x_0^\sigma)^{-1}a\,x_1^2c.
\end{equation}

\begin{definition}\label{def:GammaR}
Let \(F\) be the free profinite group on
\(\sigma,\tau,x_0,x_1\).  A finite quotient of \(F\) is
\emph{\(R\)-admissible} if the distinguished tuple generates the quotient,
the two words in \cref{eq:relators} vanish, and the normal closure of the
images of \(x_0,x_1\) is a \(2\)-group.  Let \(\GammaR\) be the inverse limit
of all \(R\)-admissible finite quotients.  Equivalently, \(\GammaR\) is the
profinite group with the displayed generators and relators, subject to the
condition that \(\normalclosure{x_0,x_1}\) be pro-\(2\).
\end{definition}

This inverse-limit formulation is the same finite-quotient semantics used in
\cite[\S2]{RT}; in particular it makes the ``normal closure is pro-\(2\)''
condition part of the presentation.

\begin{theorem}[Replacement theorem]\label{thm:main}
Assume the local and finite-target results of \cite{RT}.  Then the proof of
\cite[Theorem~1.2]{RT} remains valid with its candidate source replaced by
\(\GammaR\).  Consequently
\[
  \boxed{\GammaR\cong\GQ}.
\]
\end{theorem}

The proof occupies the rest of the note.  We verify, in order:
\begin{enumerate}
  \item the common tame quotient and full unramified character;
  \item the fully marked maximal pro-\(2\) quotient;
  \item candidate deformation duality for every elementary
        characteristic-\(2\) coefficient module;
  \item equality of the candidate and local quadratic Gauss signs.
\end{enumerate}
These are precisely the inputs isolated in
\cite[Corollary~6.19]{RT}.  The target-side recursion then depends on no
further feature of the source word; see \cite[Theorem~8.17]{RT}.

\section{Tame quotient and unramified marking}

Let \(W_R=\normalclosure{x_0,x_1}\) in \(\GammaR\).

\begin{lemma}\label{lem:tame}
The quotient by \(W_R\) is the standard tame group
\[
  \TR=\angles{\sigma,\tau\mid \tau^\sigma=\tau^2}.
\]
Moreover,
\[
  \nuR(\sigma)=1,\qquad
  \nuR(\tau)=\nuR(x_0)=\nuR(x_1)=0
\]
defines a continuous epimorphism
\(\nuR:\GammaR\twoheadrightarrow\Z_2\), and its restriction to \(\TR\) is
the standard geometric unramified character.
\end{lemma}

\begin{proof}
In a finite quotient of
\(\angles{\sigma,\tau\mid\tau^\sigma=\tau^2}\), conjugation preserves the
order of \(\tau\).  Hence \(\tau\) and \(\tau^2\) have the same order, so
\(\tau\) has odd order.  It follows that \(\tau^{\omegaTwo}=1\).

After killing \(x_0,x_1\), the wild relation is exactly
\(\tau^{\omegaTwo}=1\), so it is redundant.  This proves the tame quotient
claim.  The displayed assignment to \(\Z_2\) respects both relators:
every factor of \(\rR\) has unramified image zero.  It is surjective because
\(\sigma\) maps to \(1\).
\end{proof}

The group \(\TR\) has no nontrivial normal pro-\(2\) subgroup by the argument
of \cite[Lemma~3.3]{RT}.  Since \(W_R\) is pro-\(2\), it follows that
\[
  W_R=O_2(\GammaR).
\]
Thus the tame half of the boundary in \cite[\S4]{RT} is unchanged.

\section{\texorpdfstring{The maximal pro-\(2\) quotient}{The maximal pro-2 quotient}}

\subsection{The relator and its cup form}

\begin{lemma}\label{lem:pro2word}
The maximal pro-\(2\) quotient of \(\GammaR\) is
\begin{equation}\label{eq:DR}
  \DR=
  \angles{s,x,y\ \middle|\ 
  r_2=(x^s)^{-1}x^{-3}y^2[y,y^s]=1}_{\mathrm{pro}\text{-}2},
\end{equation}
where \(s=\sigma\), \(x=x_0\), and \(y=x_1\).
\end{lemma}

\begin{proof}
In a pro-\(2\) quotient, the closure of \(\angles{\tau}\) is procyclic.
Conjugation by \(\sigma\) acts on it as squaring.  Squaring is not an
automorphism of a nontrivial procyclic pro-\(2\) group, so \(\tau=1\).
Profinite exponentiation by \(\omegaTwo\) is the identity on pro-\(2\)
elements.  Substitution in \(\rR\) gives \cref{eq:DR}.  Conversely, every
finite quotient of \cref{eq:DR} is \(R\)-admissible, so this is the maximal
pro-\(2\) quotient.
\end{proof}

Let \(\Zassenhaus{n}\) denote the \(n\)-th Zassenhaus subgroup of the free
pro-\(2\) group on \(s,x,y\).

\begin{lemma}\label{lem:initial}
The degree-two initial form of \(r_2\) is
\[
  y^2+[x,s]\in \Zassenhaus{2}/\Zassenhaus{3}.
\]
Consequently \(\DR\) is a rank-three Demushkin group, and its cup--Bockstein
matrix in the basis dual to \(s,x,y\) is
\begin{equation}\label{eq:cupmatrix}
  \begin{pmatrix}
    0&1&0\\
    1&0&0\\
    0&0&1
  \end{pmatrix}.
\end{equation}
\end{lemma}

\begin{proof}
Since \(x^s=x[x,s]\), one has
\[
  r_2=[x,s]^{-1}x^{-4}y^2[y,y^s].
\]
Now \(x^{-4}\in\Zassenhaus{4}\).  Also \(y^s=y[y,s]\) with
\([y,s]\in\Zassenhaus{2}\), whence
\([y,y^s]\in\Zassenhaus{3}\).  This gives the asserted initial form.

The presentation is minimal and has one relation, so
\(\dim_{\F_2}H^2(\DR,\F_2)=1\).  The standard relation--cup formula gives
\cref{eq:cupmatrix}, which is nonsingular.  This is exactly the defining
cohomological condition for a Demushkin group; compare
\cite[\S1]{Labute}.
\end{proof}

Abelianizing the relation and putting \(t=\bar y-2\bar x\) gives
\begin{align}
  B_R:=\DR^\ab
    &=\angles{\bar s,\bar x,\bar y\mid
       -4\bar x+2\bar y=0}_{\Z_2},\label{eq:BR}\\
  t&=\bar y-2\bar x,\label{eq:tR}\\
  B_R&=C_2t\oplus\Z_2\bar s\oplus\Z_2\bar x,\qquad
  \nuR(t,\bar s,\bar x)=(0,1,0).\label{eq:BRsplit}
\end{align}

\subsection{The canonical orientation}

Agreement as an unmarked Demushkin group is insufficient: the tame and
pro-\(2\) quotients must have the same map to the full quotient \(\Z_2\).
We therefore compute the canonical Demushkin orientation integrally.

\begin{proposition}\label{prop:orientation}
Let \(\chiR:\DR\to\Z_2^\times\) be the canonical Demushkin orientation and
write
\[
  S=\chiR(s),\qquad X=\chiR(x),\qquad Y=\chiR(y).
\]
Then
\begin{equation}\label{eq:orientationvalues}
  Y=-X^2,\qquad
  X^3+2X^2+1=0,\qquad
  S=-\frac{X^3}{X^2+X+1}.
\end{equation}
The polynomial \(Z^3+2Z^2+1\) has a unique root in \(\Z_2^\times\), and
\[
  X\equiv5\pmod {16},\qquad S\equiv13\pmod {16}.
\]
In particular,
\[
  \im\chiR=\{\pm1\}\times(1+4\Z_2).
\]
\end{proposition}

\begin{proof}
We apply Labute's characterization of the canonical orientation
\cite[Theorem~4]{Labute}.  For a character \(\chi\) of the free pro-\(2\)
group, let \(D\) be a crossed derivation into the rank-one module defined by
\(\chi\):
\[
  D(gh)=Dg+\chi(g)Dh,\qquad
  D(g^{-1})=-\chi(g)^{-1}Dg.
\]
Such derivations descend through the relation precisely when \(D(r_2)=0\)
for arbitrary values of \(Ds,Dx,Dy\).

First,
\[
  D(x^s)=S^{-1}Dx+S^{-1}(X-1)Ds.
\]
For a commutator of elements having character values \(G,H\),
\begin{equation}\label{eq:commderivative}
  D[g,h]=G^{-1}(H^{-1}-1)Dg+
          H^{-1}(1-G^{-1})Dh.
\end{equation}
Using \cref{eq:commderivative} with \(g=y\), \(h=y^s\), and collecting the
three free coefficients in
\[
  r_2=(x^s)^{-1}x^{-3}y^2[y,y^s]
\]
gives
\begin{align}
  Y^2&=X^4, \label{eq:charrelation}\\
  0&=-(X^{-1}S^{-1}+X^{-2}+X^{-3}+X^{-4}),\label{eq:Cx}\\
  0&=S^{-1}\left[-X^{-1}(X-1)+
       \frac{(Y-1)^2}{Y^2}\right],\label{eq:Cs}\\
  0&=X^{-4}(1+Y)+
       Y^{-1}(1-Y^{-1})(S^{-1}-1).\label{eq:Cy}
\end{align}
For completeness, these are respectively the character relation and the
coefficients of \(Dx,Ds,Dy\).

By \cref{eq:charrelation}, \(Y=\pm X^2\), since the only square roots of
\(1\) in \(\Z_2^\times\) are \(\pm1\).  Equation~\eqref{eq:Cx} gives
\begin{equation}\label{eq:SfromX}
  S=-\frac{X^3}{X^2+X+1};
\end{equation}
the denominator is a unit.

Suppose first that \(Y=X^2\).  After clearing denominators,
\cref{eq:Cs} becomes
\[
  (X-1)(X^2-X-1)=0.
\]
The second factor is odd and therefore cannot vanish in \(\Z_2\), so
\(X=1\).  Equation~\eqref{eq:Cy} then has left side \(2\), a contradiction.
Thus this branch is impossible.

For \(Y=-X^2\), equation~\eqref{eq:Cs} becomes
\[
  X^3+2X^2+1=0.
\]
After \cref{eq:SfromX} is substituted, the left side of \cref{eq:Cy} is
\[
  \frac{(X+1)(X^3+2X^2+1)}{X^7},
\]
so \cref{eq:Cy} is automatic.  This proves
\cref{eq:orientationvalues}.

Every \(2\)-adic unit is congruent to \(1\) modulo \(2\).  For
\(f(Z)=Z^3+2Z^2+1\), one has
\[
  f(1)\equiv0\pmod2,\qquad f'(1)=7\equiv1\pmod2.
\]
Hensel's lemma therefore gives a unique unit root.  Lifting to modulus
\(16\) gives \(X\equiv5\), and \cref{eq:SfromX} gives
\(S\equiv13\).  Hence
\[
  v_2(X-1)=v_2(S-1)=2.
\]
The \(2\)-adic logarithm identifies \(1+4\Z_2\) with \(4\Z_2\); an element
whose difference from \(1\) has valuation \(2\) is a topological generator.
Finally, from \cref{eq:tR},
\[
  \chiR(t)=YX^{-2}=-1.
\]
Thus the image of \(\chiR\) is exactly
\(\{\pm1\}\times(1+4\Z_2)\).
\end{proof}

\begin{corollary}\label{cor:abstractD0}
The group \(\DR\) is abstractly isomorphic to
\begin{equation}\label{eq:D0}
  \Dzero=
  \angles{A,S_0,Y_0\mid A^2S_0^4[S_0,Y_0]=1}_{\mathrm{pro}\text{-}2},
\end{equation}
and hence to \(\GQpro\).
\end{corollary}

\begin{proof}
The abelianization \cref{eq:BR} has torsion subgroup of order \(2\), so the
Labute invariant is \(q=2\).  The group has odd rank \(3\), and
\cref{prop:orientation} shows that the secondary orientation depth is
\(f=2\).  Labute's odd-rank classification
\cite[Theorems~4 and~8]{Labute} therefore gives the normal form
\cref{eq:D0}.  The same normal form is the maximal pro-\(2\) Galois group of
\(\Q_2\); see \cite[Lemma~3.4]{RT}.
\end{proof}

\subsection{Matching the full unramified character}

Let
\[
  t_0=\bar A+2\bar S_0,\qquad \eta=(-3)^{-1}.
\]
The standard abelianization and its two characters are
\begin{align}
  \Dzero^\ab
   &=C_2t_0\oplus\Z_2\bar S_0\oplus\Z_2\bar Y_0,\label{eq:B0}\\
  \nu_0(t_0,\bar S_0,\bar Y_0)&=(0,1,0),\label{eq:nu0}\\
  \chiZero(t_0,\bar S_0,\bar Y_0)&=(-1,1,\eta).\label{eq:chi0}
\end{align}
These are \cite[(3.4)--(3.6)]{RT}.

\begin{proposition}\label{prop:markedpro2}
There is an isomorphism of fully marked pro-\(2\) groups
\[
  (\DR,\nuR)\cong(\GQpro,\nu_{\ur}).
\]
\end{proposition}

\begin{proof}
Since \(X,S\) both topologically generate \(1+4\Z_2\), there is a unique
\(b\in\Z_2^\times\) such that \(S=X^b\).  In the additive notation of
\(B_R\), put
\[
  k=\bar s-b\bar x.
\]
\Cref{eq:BRsplit,eq:orientationvalues} give
\[
  B_R=C_2t\oplus\Z_2k\oplus\Z_2\bar x,\qquad
  \nuR(t,k,\bar x)=(0,1,0),\qquad
  \chiR(t,k,\bar x)=(-1,1,X).
\]
The element \(\eta=(-3)^{-1}\) also topologically generates
\(1+4\Z_2\), so there is a unique unit \(u\in\Z_2^\times\) with
\(\eta^u=X\).  Therefore
\begin{equation}\label{eq:desiredab}
  \phi_{\ab}:B_R\xrightarrow{\sim}\Dzero^\ab,\qquad
  t\longmapsto t_0,\quad
  k\longmapsto\bar S_0,\quad
  \bar x\longmapsto u\bar Y_0
\end{equation}
preserves both \(\nu\) and the canonical orientation.

Choose an abstract isomorphism \(f:\DR\to\Dzero\) from
\cref{cor:abstractD0}.  Canonical orientations are functorial, so both
\(f_{\ab}\) and \(\phi_{\ab}\) preserve orientation.  Hence
\(\phi_{\ab}f_{\ab}^{-1}\) is an orientation-preserving automorphism of
\(\Dzero^\ab\).  By \cite[Proposition~3.9]{RT}, every such automorphism
lifts to an automorphism \(\gamma\) of \(\Dzero\).  Replacing \(f\) by
\(\gamma f\) gives an isomorphism whose abelianization is
\(\phi_{\ab}\), and therefore preserves the full \(\Z_2\)-valued
unramified character.  Finally use the marked local normalization
\cite[Proposition~1.1]{RT}.
\end{proof}

Combining \cref{lem:tame,prop:markedpro2}, the tame and pro-\(2\) boundary
\[
  \TR\times_{\Z_2}\GQpro
\]
is the same for \(\GammaR\) and \(\GQ\).  This verifies the structural input
to \cite[\S4]{RT}.

\section{The evaluated Fox complex}

Fix a finite boundary-compatible lower quotient \(C\), a lower map
\(\rho:\GammaR\twoheadrightarrow C\), and an elementary
\(\F_2[C]\)-module \(V\).  On a simple composition factor, the wild
normal subgroup acts trivially by \cite[Lemma~5.12]{RT}; hence the action is
tame.  Write
\[
  \mathsf S=\rho(\sigma)|_V,\qquad
  \mathsf T=\rho(\tau)|_V,\qquad
  \mathsf U=\mathsf S^{\omegaTwo}|_V.
\]
If \(e\) is the odd order of \(\mathsf T\), put
\[
  P=1+\mathsf T+\cdots+\mathsf T^{e-1}.
\]
This is the projector onto \(V^{\mathsf T}\).  Write variations of
\((\sigma,\tau,x_0,x_1)\) as \((a,b,c,d)\in V^4\).

\begin{proposition}[Evaluated Jacobian]\label{prop:jacobian}
On a simple tame module the differential obtained by differentiating the
two relators is
\begin{equation}\label{eq:jacobian}
  d^1_{R,\rho,V}=
  \begin{pmatrix}
    \mathsf S^{-1}(1+\mathsf T)&
    \mathsf S^{-1}+1+\mathsf T&0&0\\
    0&P&P+\mathsf S^{-1}&0
  \end{pmatrix}.
\end{equation}
Thus it is exactly the matrix in \cite[(5.4)]{RT} after interchanging the
two wild columns.
\end{proposition}

\begin{proof}
The tame row is unchanged from \cite[(5.5)]{RT}:
\[
  L_{\mathrm t}
  =\mathsf S^{-1}(1+\mathsf T)a+
    (\mathsf S^{-1}+1+\mathsf T)b.
\]
For the wild row, the action of \(x_0,x_1\) on \(V\) is trivial.
Consequently
\[
  D\bigl((x_0^\sigma)^{-1}\bigr)=\mathsf S^{-1}c.
\]
For \(g=x_0^{-3}\tau\), one has
\[
  Dg=(-3)c+b=c+b
\]
in characteristic \(2\), and \(g\) acts on \(V\) as \(\mathsf T\).
The profinite power/norm rule used in \cite[Lemma~5.5]{RT} gives
\[
  D(g^{\omegaTwo})=P(c+b).
\]
Finally,
\[
  D(x_1^2)=(1+1)d=0,
\]
and both entries of \([x_1,x_1^{\sigma_2}]\) act trivially, so its first
derivative is zero.  Therefore
\[
  L_{\mathrm w}=Pb+(P+\mathsf S^{-1})c,
\]
which proves \cref{eq:jacobian}.
\end{proof}

\begin{lemma}[Simple normal forms]\label{lem:normalforms}
Let \(V\) be nontrivial and simple.  Every class in the degree-one
cohomology of the word complex has a unique representative
\[
  (a,b,c,d)=(0,0,0,d).
\]
The Jacobian is surjective, and
\[
  H^0_R(V)=H^2_R(V)=0,\qquad
  \dim_{\F_2}H^1_R(V)=\dim_{\F_2}V.
\]
\end{lemma}

\begin{proof}
The normal subgroup generated by tame inertia makes
\(V^{\mathsf T}\) a submodule.  Simplicity gives two cases.

If \(\mathsf T=1\), then \(P=1\).  The tame and wild rows of
\cref{eq:jacobian} first give
\[
  b=0,\qquad (1+\mathsf S^{-1})c=0.
\]
Since \(V\) is nontrivial simple, \(\mathsf S-1\) is invertible; hence
\(c=0\), and the unique coboundary with
\((\mathsf S-1)v=a\) kills \(a\).

If \(V^{\mathsf T}=0\), then \(P=0\) and
\(\mathsf T-1=1+\mathsf T\) is invertible.  The wild row gives \(c=0\).
Subtracting the coboundary of
\((\mathsf T-1)^{-1}b\) kills \(b\), after which the tame row forces
\(a=0\).  Uniqueness follows in both cases from the same invertibility
statements.  Surjectivity of the two relation rows and the cohomological
dimension assertions follow immediately.
\end{proof}

\begin{lemma}[Trivial coefficient]\label{lem:trivial}
For \(V=\F_2\) with trivial action,
\[
  d^1_R(a,b,c,d)=(b,b).
\]
Thus \(H^1_R(\F_2)\) has coordinates \((a,c,d)\), and its scalar
cup--Bockstein form is
\begin{equation}\label{eq:scalarform}
  \langle(a,c,d),(a',c',d')\rangle
  =ac'+ca'+dd'.
\end{equation}
\end{lemma}

\begin{proof}
In \cref{eq:jacobian}, set
\(\mathsf S=\mathsf T=P=1\).  This gives the displayed differential.
Every homomorphism to \(\F_2\) factors through the maximal pro-\(2\)
quotient.  The initial form of \cref{lem:initial} therefore gives
\cref{eq:scalarform}.  Equivalently, it is the form with matrix
\cref{eq:cupmatrix}.
\end{proof}

\section{Fox--Heisenberg duality}

\subsection{The Stokes endpoint}

The chain map in \cite[\S5.2]{RT} is obtained by tracing the two
relation-defect coordinates to one scalar.  For a new word, the only
endpoint issue is the finite-word Stokes condition.

\begin{lemma}\label{lem:stokes}
Modulo \(2\), the exponent vectors of \(\rt\) and \(\rR\), in the ordered
basis \((\sigma,\tau,x_0,x_1)\), are both
\[
  (0,1,0,0).
\]
Their sum is zero.  Hence the trace
\[
  (u_{\mathrm t},u_{\mathrm w})\longmapsto
  u_{\mathrm t}+u_{\mathrm w}
\]
satisfies the endpoint condition of the finite-word Stokes formula
\cite[Lemma~5.7 and Proposition~5.8]{RT}.
\end{lemma}

\begin{proof}
For \(\rt=(\tau^\sigma)^{-1}\tau^2\), the \(\tau\)-exponent is
\(-1+2\equiv1\pmod2\).

Choose any finite integer exponent representing \(\omegaTwo\) at the
finite level under consideration.  It is odd.  In \(\rR\), the first
factor contributes one \(x_0\) modulo \(2\), while
\((x_0^{-3}\tau)^{\omegaTwo}\) contributes one \(x_0\) and one \(\tau\).
The two \(x_0\)-contributions cancel.  The square \(x_1^2\) and the
commutator have zero exponent vector modulo \(2\).  Thus only the
\(\tau\)-coordinate remains.
\end{proof}

\subsection{The mixed Hessian}

\begin{proposition}\label{prop:hessian}
Let \(V\) be nontrivial simple and represent a degree-one class by
\((0,0,0,d)\) as in \cref{lem:normalforms}.  Represent the dual class by
\(\lambda\in V^\vee\) in the corresponding \(x_1\)-supported normal form.
The degree-one pairing induced by the Fox--Heisenberg chain map is
\begin{equation}\label{eq:pairingoperator}
  (d,\lambda)\longmapsto
  \begin{cases}
    \lambda(d),&\mathsf T=1,\\[1mm]
    \lambda\bigl((1+\mathsf U+\mathsf U^{-1})d\bigr),
      &V^{\mathsf T}=0.
  \end{cases}
\end{equation}
Both operators are invertible.
\end{proposition}

\begin{proof}
On the normalized representatives, the factors
\((x_0^\sigma)^{-1}\) and
\((x_0^{-3}\tau)^{\omegaTwo}\) have no varying coordinate.  The square
\(x_1^2\) contributes the usual Heisenberg diagonal \(\lambda(d)\).
The class-two commutator identity applied to
\([x_1,x_1^{\sigma_2}]\) contributes
\[
  \lambda(\mathsf U d)+\lambda(\mathsf U^{-1}d).
\]
If \(\mathsf T=1\), then \(\mathsf U=1\) on a nontrivial simple module,
and these two commutator terms cancel.  If \(V^{\mathsf T}=0\), they give
the second line of \cref{eq:pairingoperator}.

In the ramified case, \(\mathsf U\) has \(2\)-power order, so its minimal
polynomial divides a power of \(Z+1\).  Since
\(Z^2+Z+1\) is coprime to \(Z+1\),
\[
  1+\mathsf U+\mathsf U^{-1}
  =\mathsf U^{-1}(\mathsf U^2+\mathsf U+1)
\]
is invertible.  The unramified operator is the identity.
\end{proof}

\begin{proposition}[Candidate deformation duality]\label{prop:duality}
For every elementary, possibly nonsplit, \(\F_2[C]\)-module \(V\), the
finite-word complex of \(\GammaR\) is naturally self-dual in the sense of
\cite[Proposition~5.15]{RT}.  In particular, it has the same obstruction
dimensions, lift multiplicities, perfect pairings, and adjoint connecting
homomorphisms as the local cochain complex of \(\GQ\).
\end{proposition}

\begin{proof}
By \cref{lem:stokes}, the finite-word Stokes construction gives a natural
chain map.  For the trivial simple module it is a quasi-isomorphism by the
nonsingularity of \cref{eq:scalarform}.  For every nontrivial simple module,
\cref{lem:normalforms,prop:hessian} show that the degree-one pairing is
perfect and that \(H^0=H^2=0\).  Hence the cone of the chain map is acyclic
on all simple modules.

The word complex, its dual, and the chain map are exact coefficient
functors.  Applying the degreewise short exact cone sequence along a
composition series and using two-out-of-three extends acyclicity to every
elementary module.  This is exactly the cone devissage argument of
\cite[Lemma~5.11 and Proposition~5.15]{RT}.  Naturality of the chain map
gives adjointness of connecting morphisms.  Comparison with the local
complex then uses local Tate duality as in
\cite[Proposition~5.16 and Corollary~5.17]{RT}.
\end{proof}

\section{Quadratic determinant obstructions}

Let \(V\) be a nontrivial self-dual simple module, and let
\[
  q:V\longrightarrow\F_2
\]
be a nonzero invariant nonsingular quadratic form with polar form \(b_q\).
Let \(\kappa_q^0\) be the normalized base determinant class of
\cite[Lemma~6.3]{RT}.

\begin{proposition}[Base word expansion]\label{prop:quadratic}
On the normalized class \((0,0,0,d)\), evaluation of Roe's wild word gives
\begin{equation}\label{eq:QR}
  Q_R^0(d)=
  \begin{cases}
    q(d),&\mathsf T=1,\\[1mm]
    q(d)+b_q(d,\mathsf U^{-1}d),&V^{\mathsf T}=0.
  \end{cases}
\end{equation}
In the ramified case its polar form is
\begin{equation}\label{eq:polar}
  b_R(d,d')
  =b_q\bigl(d,(1+\mathsf U+\mathsf U^{-1})d'\bigr).
\end{equation}
\end{proposition}

\begin{proof}
As in the mixed calculation, only the final two factors vary on the
normalized representative.  The square \(x_1^2\) contributes \(q(d)\).
The complete class-two commutator calculation for
\([x_1,x_1^{\sigma_2}]\) contributes
\[
  b_q(d,\mathsf U^{-1}d).
\]
No equivariance of a chosen section is being assumed here: changing the
section alters the factor set by a normalized coboundary, and the complete
lifted multiplication formula of \cite[(6.8)]{RT} yields the same
cohomology class.

If \(\mathsf T=1\), then \(\mathsf U=1\) and
\(b_q(d,d)=0\), giving the first line.  Polarizing the second line and
using \(\mathsf U\)-invariance and symmetry of \(b_q\) gives
\cref{eq:polar}.
\end{proof}

The formula in \cref{eq:QR} is exactly \cite[(6.25)]{RT}, with the
normalized wild coordinate renamed from \(c\) to \(d\).  Consequently all
of the candidate-side Arf computations in \cite[\S6.2.1]{RT} apply without
change.

\begin{corollary}[Gauss signs]\label{cor:gauss}
If \(n=\dim_{\F_2}V\), then
\[
  \#(Q_R^0)^{-1}(0)=
  \begin{cases}
    2^{n-1}-2^{n/2-1},&V\text{ unramified},\\
    2^{n-1}+2^{n/2-1},&V\text{ ramified}.
  \end{cases}
\]
These are the same zero counts, and hence the same Gauss signs, as the
local determinant obstruction for \(\GQ\).
\end{corollary}

\begin{proof}
For unramified \(V\), \cref{eq:QR} is the invariant quadratic form \(q\);
the Hermitian-line calculation of \cite[Proposition~6.9]{RT} gives the
negative sign.  In the ramified case, \cref{eq:QR} is the Wall-twisted form
used there, so \cite[Lemmas~6.6 and~6.8]{RT} give Arf invariant zero and
the positive sign.  Equality with the local signs is
\cite[Proposition~6.18]{RT}.
\end{proof}

For a general determinant class
\[
  \kappa=\kappa_q^0+\Gamma_{\gamma_\kappa}+\operatorname{inf}\delta_\kappa,
\]
the difference from the base form is the natural linear pairing term plus
the inflated scalar term, exactly as in \cite[(6.7) and (6.27)]{RT}.
This assertion is formal in the determinant class and the natural chain
pairing; no additional word calculation is required.  Thus the affine
transgression and shear corrections of \cite[\S6.5]{RT} also carry over.

\section{Passage to every finite target}

\begin{proposition}\label{prop:interface}
For every finite lower quotient, every boundary-compatible exact-image map,
and every elementary characteristic-\(2\) coefficient module, the sources
\(\GammaR\) and \(\GQ\) have:
\begin{enumerate}
  \item the same obstruction-space dimensions and the same number of lifts
        when the obstruction vanishes;
  \item natural perfect pairings for all short exact coefficient
        sequences, with adjoint connecting maps;
  \item the same scalar cup--Bockstein form;
  \item for every invariant nonsingular quadratic form on a self-dual
        simple head, the same base determinant Gauss sign.
\end{enumerate}
\end{proposition}

\begin{proof}
Items (1) and (2) are \cref{prop:duality}; item (3) is
\cref{lem:trivial}; and item (4) is \cref{cor:gauss}.  The marked maximal
pro-\(2\) input is \cref{prop:markedpro2}, and the common tame--pro-\(2\)
boundary follows from \cref{lem:tame}.  This is the analogue for
\(\GammaR\) of the complete source interface
\cite[Corollary~6.19]{RT}.
\end{proof}

\begin{proof}[Proof of \cref{thm:main}]
The target-side operations in \cite[\S\S7--9]{RT} are restrictions,
pushouts, exact-image partitions, finite Fourier transforms, and central
cover constructions internal to the finite target.  By
\cite[Theorem~8.17]{RT}, the exact-image recursion depends on the profinite
source only through the common boundary and the source interface listed in
\cref{prop:interface}.  Replacing the original candidate source by
\(\GammaR\) therefore leaves the recursion unchanged.

The boundary-framed induction of \cite[Theorem~4.2 and \S9]{RT} gives equal
exact-image counts for \(\GammaR\) and \(\GQ\).  Removing the framing as in
\cite[\S10]{RT} yields
\[
  |\Sur(\GammaR,H)|=|\Sur(\GQ,H)|
\]
for every finite group \(H\).  The group \(\GammaR\) is topologically
finitely generated.  The one-sided profinite reconstruction lemma
\cite[Lemma~2.5]{RT} now gives \(\GammaR\cong\GQ\).
\end{proof}

\clearpage
\section{Audit summary and remaining formalization work}

The proof above checks all features of Roe's relator that are visible to
the source interface:
\begin{itemize}
  \item its tame specialization is redundant;
  \item its maximal pro-\(2\) quotient has the correct Demushkin form,
        orientation image, and full \(\Z_2\)-marking;
  \item its evaluated Fox row is the paper's row with the two wild columns
        interchanged;
  \item its mod-\(2\) exponent vector satisfies the Stokes endpoint;
  \item its mixed Heisenberg Hessian and determinant quadratic form are
        exactly those used in the paper after the same wild-coordinate
        interchange.
\end{itemize}

This is not an explicit Nielsen or Tietze transformation between Roe's
presentation and the longer presentation in \cite[Theorem~1.2]{RT}.
No such transformation is needed for the finite-quotient proof.  The result
instead shows that the two sources have the same exact-image counts against
every finite target, from which abstract profinite isomorphism follows.

For a machine-checked version, the genuinely new formal obligations are
short and well isolated:
\begin{enumerate}
  \item formalize \cref{eq:charrelation,eq:Cx,eq:Cs,eq:Cy} and
        the Hensel root calculation;
  \item formalize the wild Fox row in \cref{eq:jacobian};
  \item formalize the \(x_1^2[x_1,x_1^{\sigma_2}]\) mixed and quadratic
        Heisenberg expansions;
  \item instantiate the generic finite-word Stokes theorem using
        \cref{lem:stokes}.
\end{enumerate}
The local calculations and finite-target induction need no new
candidate-specific mathematics.

\begin{thebibliography}{9}

\bibitem{RT}
D. Roe and D. Turturean,
\emph{A Presentation of the Absolute Galois Group of \(\Q_2\)},
2026.
\newblock
\url{https://roed314.github.io/gq2/paper/paper.html}.

\bibitem{Labute}
J. P. Labute,
\emph{Classification of Demushkin Groups},
Canadian Journal of Mathematics \textbf{19} (1967), 106--132.
\newblock
\url{https://doi.org/10.4153/CJM-1967-007-8}.

\bibitem{Presentations}
D. Roe,
\emph{Presentations of \(G_{\Q_2}\): David Roe's candidate},
2026.
\newblock
\url{https://roed314.github.io/gq2/presentations/#candidate-roe}.

\end{thebibliography}

\end{document}
